\section{问题分析}

四问共享同一套微网物理结构：小区负荷必须被满足，光伏优先供当前负荷，储能承担跨时段能量转移，
外网购电补足缺口。真正随问题变化的是\textbf{决策时刻可获得的信息量}与\textbf{购电结算规则}，
因此不宜把四问建成彼此独立的模型，而应从问题一的显式能量分流模型出发逐层扩展。

\subsection{问题一分析}

问题一中当天电价、负荷与光伏曲线均已知，属单日确定性优化。电价有明显峰谷差，光伏集中在白天，
且部分时段光伏超过负荷，储能的直接作用因此有两个：吸收富余光伏，以及在低价时段购电充电、
高价时段放电。为刻画外网供负荷、外网充电、光伏充电、储能放电与弃光各自的去向，
需把总能量平衡拆成显式流量平衡；又因同一时段储能不能同时充放，需引入 0-1 状态变量，
故问题一适合建立\textbf{能量分流混合整数线性规划模型}，以全天购电费用最小为目标。
问题一的意义不止于求出该日策略，更在于形成后续三问可直接复用的物理层。

\subsection{问题二分析}

问题二的关键变化是每天 0:00 制定计划时，当天真实负荷与光伏尚未全部实现，
而计划购电不足会触发 5 倍电价的紧急购电。因此不能在 0:00 把储能充放电等执行变量全部锁死：
真正需要提前确定的只是\textbf{计划购电量}，充放电、弃光与紧急购电应在情景实现后再定，
由此形成“0:00 定计划—情景实现后物理补救”的两层决策。为避免单日模型在时域末端过度放电，
还需让日末储电量具有未来价值，故采用 7 日滚动视野，并用历史预测残差构造 SAA 情景，
把预测误差风险直接纳入日前计划，建立为\textbf{7 日滚动 SAA 两阶段随机 MILP}。

\subsection{问题三分析}

问题三中微网设备与购电物理过程不变，新增的是 6:00、12:00、18:00 三次信息更新机会，
尚未执行的计划可以重新调整，因此需要明确新旧信息的衔接方式：新预报替换旧预报后，
已执行的历史时段必须锁定，未执行的区间以当前真实储电量作为阶段初值重新优化；
调增、调减与紧急购电统一相对当天 0:00 的原计划结算，使调整量不会层层叠加。
题目给出的光伏预报仍只是预测而非真值，最新预报只能作为中心预测，仍需叠加历史预报误差构造情景。
据此，问题三在问题二基础上引入阶段集合 \(\mathcal S=\{0,6,12,18\}\)，
形成\textbf{基于多时点光伏预报更新的滚动 SAA-MILP 购电调整模型}。

\subsection{问题四分析}

问题四进一步取消“未来电价已知”的假设，负荷、光伏与电价同时存在不确定性。
未来真实电价只能在其实际发生后用于结算，不能提前进入计划优化，否则构成未来信息泄漏。
对 Q4-2，需在问题二中增加未来电价预测，并把价格误差与负荷、光伏误差按同一历史日联合抽样，
形成 \((L^{(\omega)},PV^{(\omega)},\pi^{(\omega)})\) 联合情景；对 Q4-3，还需在问题三的
0:00、6:00、12:00、18:00 四阶段结构中，用已实现的真实价格对剩余时段价格预测作因果更新。

历史价格误差叠加后可能出现负价格。若仍允许无上界的计划购电，模型会为负成本而购买无实际用途的电量；
若只用连续变量表示调增调减，负价格还可能导致二者同时为正。故问题四另需两项针对性约束：
H-1 计划购电物理上界，与 H-2 调增调减显式互斥。

\subsection{四问统一建模关系}

四问建立在同一物理系统之上，逐层递进（总体流程见图~\ref{fig:flow}）：
问题一建立确定性能量分流 MILP；问题二把计划变量与执行变量分层，引入 SAA 情景与 7 日滚动；
问题三加入四个时点的信息更新与计划再调整；问题四再纳入价格预测、联合价格情景与负价加固。
全文统一复用问题一建立的物理流量变量，只在后续问题中按新增业务规则增加必要变量，
以保证模型之间的连续性、符号的一致性与结果的可解释性。

% 总体流程图（建模手提供）已作为 fig:flow 插在「问题背景与问题重述」末尾；
% 本节的四问关系改用文字表述，不再另画流程框图。
